\documentclass[11pt]{article}
\usepackage[margin=0.9in]{geometry}
\usepackage{hyperref}
\usepackage{listings}
\usepackage{xcolor}
\usepackage{booktabs}
\usepackage{enumitem}
\usepackage{longtable}
\usepackage{array}
\usepackage{amsmath}
\usepackage{amssymb}
\usepackage{ragged2e}

\emergencystretch=4em
\tolerance=9999
\hbadness=10000
\sloppy

\lstset{
  basicstyle=\ttfamily\small,
  breaklines=true,
  columns=fullflexible,
  frame=single,
  backgroundcolor=\color{gray!10},
  showstringspaces=false
}

\newcolumntype{L}[1]{>{\RaggedRight\arraybackslash}p{#1}}

\title{Spatial Range Over Which the \texttt{gradients\_by\_period/}
       Gradients Were Calculated\\
       \large grep-dap test case, Prompt \#4}
\author{Compiled for P. Cornillon}
\date{2026-05-16}

\begin{document}
\maketitle

\begin{abstract}
The gradient sums in
\texttt{https://sst-aqua.gso.uri.edu/opendap/gradients\_by\_period/}
are 1$^\circ$ monthly accumulations of per-pixel SST gradients that
were originally computed at the L2 swath level, inside the
companion \texttt{SST\_Orbits/} files. The question for Prompt~\#4
is the spatial extent over which each per-pixel gradient was
computed --- i.e.\ the support of the differentiation operator.
Three independent estimates (pixel-count ratios, magnitude-matching
of stored vs.\ centered-finite-difference candidates, and
autocorrelation of the stored gradient field) converge on
\textbf{$5 \pm 2$ km}. The likely operator is a centered
finite difference at $\sim\pm 2$ to $\pm 3$ pixels applied to the
masked SST field \texttt{regridded\_sst} of the orbit files
(which is itself derived from raw \texttt{SST\_In} by masking out
pixels flagged as high-gradient noise). The estimate is robust to
the exact operator shape; what it cannot resolve is whether the
\texttt{regridded\_sst} construction includes additional spatial
smoothing on top of the masking step.
\end{abstract}

\tableofcontents

\section{What ``spatial range over which the gradients were
calculated'' means}\label{sec:question}

The variables in \texttt{gradients\_by\_period/} are 1$^\circ$
monthly bin sums of the per-pixel orbit-level gradient values
\texttt{eastward\_gradient} and \texttt{northward\_gradient} from
\texttt{SST\_Orbits/} (and their swath-relative counterparts
\texttt{grad\_as\_per\_km} / \texttt{grad\_at\_per\_km} which are
not exposed in the orbit file but follow the same algorithm). The
1$^\circ$ binning is a \emph{post-hoc} accumulation; the spatial
range over which a single gradient value was computed is set
inside the orbit-level processing, by the differentiation operator
that maps a window of SST pixels to one gradient value. That
support, in km, is the answer the prompt asks for.

\section{Three lines of evidence}\label{sec:evidence}

\subsection{Line 1: pixel-count ratios}\label{sec:line1}

From Prompt~\#3, the three pixel-count flavours in the monthly
file are different by construction:
\begin{center}
\begin{tabular}{@{}lr@{}}
\toprule
Field & Mean per cell (day side) \\
\midrule
\texttt{day\_pixel\_count}     & 10\,030 \\
\texttt{day\_as\_pixel\_count} & ~8\,110 \\
\texttt{day\_at\_pixel\_count} & ~4\,090 \\
\bottomrule
\end{tabular}
\end{center}
The fractions
$N_{\rm as}/N_{\rm pix} \approx 0.81$ and
$N_{\rm at}/N_{\rm pix} \approx 0.41$. MODIS Aqua acquires SST in
groups of 10 along-track lines per mirror scan; geolocation /
calibration jumps between scans, so along-track derivatives whose
stencil straddles a scan boundary are masked. For a centered
finite difference with along-track half-width $H$ pixels, the
fraction of SST-valid pixels with a valid along-track gradient is
approximately $(10-2H)/10$ minus a cloud-edge loss.
The cloud-edge loss can be read off from the along-scan ratio
(which is not affected by scan boundaries),
$1 - N_{\rm as}/N_{\rm pix} \approx 0.19$. The additional
along-track loss is then $0.81 - 0.41 = 0.40 = 2H/10$, giving
$\mathbf{H = 2}$ pixels. With along-track pixel spacing
$\sim$1.13 km near nadir, the implied along-track operator support
is $2H \times 1.13 \approx \mathbf{4.5}$ km.

\subsection{Line 2: magnitude-matching to centered differences}\label{sec:line2}

\texttt{scripts/grad\_scale\_findregion.py} located the most
gradient-populated 1000-line window of orbit 115220
(2024-01-01 01:15:58 UTC); the winning window is a 200-line
patch in the central Pacific at $\sim$20\,N, 180\,E with
$\sim$51\% valid gradient pixels. The full $200\times 1354$
block of \texttt{SST\_In}, \texttt{regridded\_sst},
\texttt{eastward\_gradient},
\texttt{northward\_gradient},
\texttt{latitude}, \texttt{longitude} was pulled to
\texttt{/tmp/grad\_block2.pkl}.

Define $|\nabla T|_{\rm stored} = \sqrt{
(\mathtt{eastward\_gradient})^2 +
(\mathtt{northward\_gradient})^2}$ and the candidate
\[
|\nabla T|_{\rm cand}(h) = \sqrt{
\Big(\frac{T_{i,j+h}-T_{i,j-h}}{2h\,\bar{\Delta x}}\Big)^{\!2}
+
\Big(\frac{T_{i+h,j}-T_{i-h,j}}{2h\,\bar{\Delta y}}\Big)^{\!2}}
\]
for $T$ equal to either \texttt{SST\_In} or
\texttt{regridded\_sst}. Magnitude (rather than signed components)
sidesteps the swath-vs-geographic rotation. The empirical
$\bar{\Delta x}$ and $\bar{\Delta y}$ used (median across the
block) are 1.35 km and 1.13 km respectively, computed by
haversine on \texttt{latitude}/\texttt{longitude}.

Results --- per-pixel RMSE between
$|\nabla T|_{\rm stored}$ and $|\nabla T|_{\rm cand}(h)$, central
swath, restricted to pixels where both are valid:

\begin{center}
\begin{tabular}{@{}rrr|rrr@{}}
\toprule
& \multicolumn{2}{c|}{$T = $ \texttt{SST\_In}} & \multicolumn{2}{c}{$T = $ \texttt{regridded\_sst}} \\
\midrule
$h$ & support (km) & RMSE & $h$ & support (km) & RMSE \\
\midrule
1 & 2  & 0.079 & 1 & 2  & 0.059 \\
2 & 4  & 0.034 & 2 & 4  & 0.025 \\
\textbf{3} & \textbf{6}  & \textbf{0.028} & \textbf{3} & \textbf{6}  & \textbf{0.025} \\
4 & 8  & 0.029 & 4 & 8  & 0.025 \\
5 & 10 & 0.030 & 5 & 10 & 0.027 \\
7 & 14 & 0.032 & 7 & 14 & 0.028 \\
10 & 20 & 0.034 & 10 & 20 & 0.030 \\
\bottomrule
\end{tabular}
\end{center}

The RMSE minimum is at $h = 3$ pixels, i.e.\ a centered-difference
support of \textbf{6 km}, when the candidate is built from
\texttt{regridded\_sst}; using \texttt{SST\_In} the minimum sits at
the same $h$ with a slightly higher residual --- consistent with
\texttt{regridded\_sst} being the operative input field.

A second, looser check is the ratio of spatially-averaged
magnitudes. The stored mean $|\nabla T|$ in the central window is
$0.0344\,^\circ$C/km. Mean
$|\nabla T|_{\rm cand}(h)$ from \texttt{regridded\_sst} matches
1:1 at $h \approx 2.5$ pixels (interpolated), i.e.\ support
\textbf{$\sim$4--6 km}. Mean from \texttt{SST\_In} matches 1:1
only at $h \approx 10$ pixels (support $\sim$20 km), a reminder
that \texttt{SST\_In} carries pixel-scale noise that the operator
must remove --- which the
\texttt{refined\_mask} / \texttt{regridded\_sst} step
accomplishes.

\subsection{Line 3: autocorrelation of the stored gradient}\label{sec:line3}

For a gradient field computed with operator support $L$, adjacent
pixel pairs share input pixels and so the gradient field is
autocorrelated over scales $\lesssim L$. Beyond $L$ the
autocorrelation should drop sharply (apart from real
oceanographic structure, which is broad compared to the operator
support). \texttt{scripts/grad\_scale\_autocorr.py} computes
$\rho(\ell)$ of $|\nabla T|_{\rm stored}$ in the central swath:

\begin{center}
\begin{tabular}{@{}rrr|rrr@{}}
\toprule
\multicolumn{3}{c|}{across-track ($\Delta x = 1.35$ km)} &
\multicolumn{3}{c}{along-track ($\Delta y = 1.13$ km)} \\
\midrule
lag (px) & dist (km) & $\rho$ & lag (px) & dist (km) & $\rho$ \\
\midrule
0 & 0   & 1.00 & 0 & 0   & 1.00 \\
1 & 1.35 & 0.64 & 1 & 1.13 & 0.50 \\
2 & 2.70 & 0.32 & 2 & 2.26 & 0.18 \\
3 & 4.05 & 0.19 & 3 & 3.39 & 0.13 \\
4 & 5.40 & 0.13 & 4 & 4.52 & 0.07 \\
5 & 6.75 & 0.10 & 5 & 5.65 & 0.05 \\
6 & 8.10 & 0.08 & 6 & 6.78 & 0.03 \\
8 & 10.8 & 0.06 & 8 & 9.0 & 0.02 \\
10 & 13.5 & 0.05 & 10 & 11.3 & 0.01 \\
\bottomrule
\end{tabular}
\end{center}

For comparison, the autocorrelation of a 1-pixel centered diff of
the raw \texttt{SST\_In} drops to \emph{negative} values
($\rho \approx -0.19$ at 2.7 km, $\rho \approx -0.23$ at 4 km)
and returns to zero only past 8 km. That anti-correlation is the
Nyquist-like fingerprint of a sharp 1-pixel derivative applied to
a noisy field. The stored gradient does \emph{not} have it, ruling
out a simple 1-pixel centered diff on \texttt{SST\_In}.

The decorrelation length (drop below $1/e \approx 0.37$) is
$\sim$2 pixels across-track and $\sim$1.5 pixels along-track ---
i.e.\ \textbf{2--3 km in either direction}. The autocorrelation
also has a long flat tail near 0.05--0.10 in across-track that
likely reflects real oceanographic structure in the 1$^\circ$
mesoscale band; that tail does not constrain the operator support.

\section{Convergent estimate}\label{sec:answer}

The three lines of evidence agree on a single number to within a
factor of $\sim$2.

\begin{center}
\begin{tabular}{@{}lr@{}}
\toprule
Line of evidence & Implied operator support \\
\midrule
Pixel-count ratios (\S\ref{sec:line1})           & $\sim$4 km (along-track) \\
Magnitude-matching RMSE (\S\ref{sec:line2})      & 4--6 km (both directions) \\
Magnitude-matching mean ratio (\S\ref{sec:line2})& 4--6 km \\
Autocorrelation decorrelation length (\S\ref{sec:line3}) & 2--4 km \\
\bottomrule
\end{tabular}
\end{center}

\smallskip
\noindent\textbf{Best estimate: $5 \pm 2$ km, with the
along-track support possibly $\sim$1 km smaller than the
across-track support owing to MODIS pixel-spacing anisotropy.}

The simplest description of the operator that is consistent with
every line of evidence is: a centered finite difference with
half-width $h = 2$--$3$ pixels (full operator support 4--6 km)
applied to \texttt{regridded\_sst}, where
\texttt{regridded\_sst} is the orbit's
\texttt{SST\_In} field after the high-gradient pixels flagged in
\texttt{refined\_mask} have been removed.

\section{Things the analysis cannot distinguish}\label{sec:caveats}

\begin{itemize}
  \item Whether the actual operator is a centered finite
    difference, a Sobel kernel, or a local least-squares fit of
    similar effective support. Any of these would give roughly the
    same fingerprints in the three lines of evidence above. The
    estimate is for the \emph{effective spatial support}, not for
    the operator's exact functional form.
  \item Whether the construction of
    \texttt{regridded\_sst} from \texttt{SST\_In} includes a
    spatial averaging step on top of the masking. The values
    of \texttt{regridded\_sst} sampled here have a standard
    deviation $\sim$3$\times$ smaller than \texttt{SST\_In}, but
    almost all of that reduction can be explained by the 33\,\%
    masking carried in \texttt{refined\_mask} (the surviving
    pixels are in coherent patches of SST). If there is in
    addition a 3$\times$3 or 5$\times$5 box smoothing
    --- which would be invisible to an external analysis ---
    the \emph{total} area of raw SST contributing to a single
    gradient value would be enlarged from the 4--6 km estimated
    above to perhaps 7--10 km.
  \item The sample region is open ocean
    ($\sim$20\,N, 180\,E) without strong fronts. In a region
    with a sharp front (e.g.\ Gulf Stream, Kuroshio) the
    front itself contributes structure at scales smaller than
    the operator support, and the autocorrelation argument
    weakens. A region with no fronts (oligotrophic gyre) would
    give a tighter constraint but I did not sample such a
    region.
  \item Only the orbit-file gradients
    \texttt{eastward\_gradient}, \texttt{northward\_gradient}
    are exposed publicly. The swath-relative
    \texttt{grad\_as\_per\_km}, \texttt{grad\_at\_per\_km}
    flavour appears only in the monthly file; the analysis
    assumes the same operator is used for both flavours, which
    the naming suggests but does not prove.
\end{itemize}

\section{Implication for downstream metadata}\label{sec:impl}

For the next prompt (catalog metadata for
\texttt{gradients\_by\_period}), the spatial resolution of the
underlying gradient measurement should be advertised as the
operator support rather than the 1$^\circ$ bin size of the monthly
file. Suggested form:

\begin{quote}\ttfamily
spatial\_resolution = "approximately 5 km (per-pixel SST gradient
operator support); summed into 1$^\circ$ monthly cells."
\end{quote}

ACDD-style alternative:

\begin{quote}\ttfamily
geospatial\_lat\_resolution = "1 degree (bin size)"\\
geospatial\_lon\_resolution = "1 degree (bin size)"\\
comment = "Each gradient sum aggregates per-pixel L2 SST gradients
whose underlying differentiation-operator support is approximately
5 km."
\end{quote}

\end{document}
